\chapter{Introducción}
%\addcontentsline{toc}{chapter}{Introducción}
\label{introduccion}

Desde el descubrimiento fortuito de los rayos X por Wilhelm C. Röntgen a finales del siglo
XIX, la capacidad de observar el interior de la materia de forma no invasiva ha sido uno
de los pilares de la ciencia moderna. Lo que comenzó como una herramienta de diagnóstico
médico ha evolucionado, gracias al desarrollo de la computación y la microelectrónica,
en técnicas de altísima precisión como la microtomografía computarizada ($\mu$CT).
Esta técnica permite hoy reconstruir volúmenes tridimensionales con resoluciones en la
escala micrométrica, convirtiéndose en un recurso indispensable no solo en medicina,
sino también en la ciencia de materiales, la geología y la industria. \cite{Stock2019}

Su utilidad evidente se refleja en la gran cantidad de avances tecnológicos y mejoras
técnicas producidas cada año, propiciadas por empresas que buscan comercializar sistemas de
$\mu$CT cada vez más accesibles y de mejor rendimiento, abocados a tareas específicas,
ya sea en medicina, industria o investigación. En el caso de esta última, sin embargo,
la implementación de un sistema de microtomografía de diseño propio o no comercial presenta
ventajas sustanciales frente a las soluciones cerradas de mercado. Mientras que los
equipos comerciales suelen operar como ``cajas negras'' con parámetros fijos y algoritmos
de procesamiento propietarios, un sistema experimental de arquitectura abierta permite la manipulación directa
de la geometría del sistema, el intercambio y la reparación de componentes clave —como la fuente 
o el detector— y el acceso total a los datos crudos y al código base. Esta flexibilidad es 
fundamental para la investigación aplicada, ya que permite adaptar el instrumento a las 
necesidades específicas de la muestra y garantiza una transparencia absoluta en la 
trazabilidad de todos los procesos de adquisición y reconstrucción. Por último, este tipo
de sistemas se presenta como una opción más económica y accesible para laboratorios de
investigación, sin por ello sacrificar la calidad de los resultados obtenidos.

Sin embargo, el potencial de un sistema de $\mu$CT está intrínsecamente ligado a la calidad
de su calibración y a la precisión de su geometría. En sistemas no comerciales o ensamblados
en laboratorios de investigación, además del conocimiento concreto sobre cómo operarlos,
surgen desafíos particulares: desde desalineaciones 
mecánicas casi imperceptibles hasta variaciones en la respuesta de intensidad del detector.
Estos factores, si no se tratan adecuadamente, derivan en artefactos que degradan la
resolución espacial y dificultan la interpretación cualitativa y cuantitativa de los resultados.
En otras palabras, una alineación geométrica precisa es el factor más influyente en la
calidad de un microtomógrafo y en la fidelidad de sus tomografías para realizar mediciones
en la escala micrométrica.

\section{Objetivos}

El objetivo general de este trabajo es el desarrollo y la planificación de
una secuencia de métodos para preparar y optimizar un 
sistema de microtomografía computarizada de rayos X. En otras palabras, se
busca establecer un protocolo detallado de calibración geométrica
de un microtomógrafo de rayos X (en particular, uno de arquitectura abierta,
para un ensamblado e implementación óptimos). Con tal fin, se propone el estudio
teórico y experimental de los procesos físicos de interacción de la radiación
con la materia que participan en los procedimientos de adquisición de imágenes
radiográficas. A su vez, se exploran técnicas de procesamiento matemático y computacional
de imágenes digitales, dirigidas al pre-procesamiento de la tomografía.

\vspace{5mm}

Se proponen como objetivos específicos:

\begin{itemize}
    \item Estudio e implementación de técnicas de alineación geométrica de un 
    sistema de microCT.
    \item Estudio e implementación de técnicas computacionales para la determinación
    de correcciones geométricas.
    \item Estudio e implementación de métodos de medición de la resolución espacial
    efectiva del sistema de microCT.
    \item Implementación de técnicas de pre-procesamiento digital en imágenes
    radiográficas, para mejorar la tomografía resultante.
    \item Implementación de lo investigado en la mejora del actual microtomógrafo
    del laboratorio.
\end{itemize}